\subsection{A singleton-comma deletion inside one block}

\begin{proposition}[Proved]\label{CP2:A-singleton-deletion-block-cut}
Let $C$ be a category with four objects and acyclic support.
Let $p$ be a source or sink of its support, with
$|\operatorname{End}_C(p)|\leq4$, and write
$E=C\setminus\{p\}$. Suppose that an object $v$ of $E$
satisfies one of the following conditions inside $E$:
\begin{enumerate}
\item All arrows into $v$ from other objects of $E$ consist
of one arrow $q\to v$, and $\operatorname{End}_C(q)$ is trivial.
\item All arrows out of $v$ to other objects of $E$ consist
of one arrow $v\to q$, and $\operatorname{End}_C(q)$ is trivial.
\end{enumerate}
Suppose further that both corner monoids of the full
two-object category $F=E\setminus\{v\}$ have order at most
two. Then every rational degree-two cohomology product on
$BC$ is zero. In particular $BC\not\simeq\mathbb{CP}^2$.
\end{proposition}
\begin{proof}
The inclusion $j:F\hookrightarrow E$ is an actual homotopy
equivalence on classifying spaces by the singleton-comma
lemma, Lemma~\ref{CP2:A-five18-singleton-deletion}.
For completeness, in the first case $(j\downarrow x)$
has terminal object $(x,1_x)$ for every retained $x$,
whereas $(j\downarrow v)$ consists of the single object
$(q,q\to v)$ and its identity. Quillen's Theorem~A applies.
The second case is its opposite. No restriction on the
endomorphisms at the deleted object $v$ is required.

We next check $H^2(BF;\mathbb Q)=0$. Every monoid of
order at most two has semisimple rational algebra: the
two nontrivial possibilities are the cyclic group of order
two and the two-element idempotent monoid, whose rational
algebras are both isomorphic to $\mathbb Q\times\mathbb Q$.
Consequently the trivial rational modules are projective
on both sides. If the two objects of $F$ are comparable,
apply the actual orbit functors of
Lemma~\ref{CP2:A-high-orbits} first at the source and then
at the sink. They induce rational cohomology isomorphisms
to a two-object category with trivial corners and some
nonempty finite set of parallel arrows. Its nerve is a
graph, so its second cohomology vanishes. If the objects
are incomparable, $BF$ is the disjoint union of two
rationally acyclic corner spaces and the same conclusion
holds. Thus $H^2(BE;\mathbb Q)=0$ by the specified inclusion.

The full corner $D$ at $p$ also has zero second rational
cohomology by Corollary~\ref{CP2:A-five-types}. The source
or sink hypothesis gives a one-way support partition into
$D,E$, with the order of the blocks reversed in the sink
case. Lemma~\ref{CP2:A-two-block} now makes every product
of two degree-two classes zero. This argument neither
requires bounded homology nor identifies $BC$ with $BE$.
\end{proof}

\begin{remark*}
In the frozen four-object eighteen-arrow rows $1777$ and
$1778$, take $p=1$, $v=2$ and $q=0$. Within the full
deletion on $0,2,3$, there is exactly one incoming strict
arrow to $2$, from the trivial corner $0$. Deleting $2$
leaves a graph with respectively two or three parallel
arrows from $0$ to $3$. Thus both rows satisfy the proposition
even though that three-object block has no initial or
terminal object.
\end{remark*}
